\documentclass[10pt,a4paper]{amsart}
\usepackage[margin=19mm]{geometry}
\usepackage{amsmath,amssymb,mathtools}
\usepackage[scr=rsfs,cal=euler]{mathalfa}
\usepackage{xfrac}
\usepackage[unicode,hidelinks,bookmarks=false]{hyperref}
\newcommand{\ie}{i.e.\ }
\newcommand{\cf}{cf.\ }
\newcommand{\resp}{respectively\ }
\newcommand{\Subsection}{\subsection}
\providecommand{\nobreakdash}{\nobreak\hbox{-}}
\setlength{\emergencystretch}{2em}
\multlinegap0pt
\usepackage{mathtools}
\usepackage{amssymb}
\usepackage{tikz}
\usepackage{bm}
\usepackage[normalem]{ulem}
\newtheorem{lemma}{Lemma}
\newcommand{\ID}{\mathrm{I}}
\newcommand{\M}{M}
\newcommand{\OA}{\mathbf{A}}
\newcommand{\OAD}{\mathbf{\bar{A}}}
\newcommand{\abrOSCR}{\mathrm{R}}
\newcommand{\mmu}{\mu}
\newcommand{\oad}{\mathbf{\bar{a}}}
\title{(Integrability of the multispecies harmonic process}
\author{Francesco Casini, Rouven Frassek, Cristian Giardinà}
\date{Source: arXiv:2607.26262v1 (CC BY 4.0)}
\begin{document}
\maketitle

\noindent\emph{Source and licence.} (Integrability of the multispecies harmonic process. Version arXiv:2607.26262v1 (CC BY 4.0), distributed by arXiv under the Creative Commons Attribution 4.0 International licence, http://creativecommons.org/licenses/by/4.0/, as recorded in the arXiv licence metadata for this exact version and shown on the abstract page https://arxiv.org/abs/2607.26262v1. Full licence notice: \texttt{licenses/arxiv-2607.26262-CC-BY-4.0.txt}.\par\smallskip

\noindent\textit{Editorial scope.} Counted unit: the article's lemma Lemma-equivalence (source0.tex lines 885--898). The article's complete original proof of it is delivered with it (source0.tex lines 899--1244, one \texttt{proof} environment of 346 lines). No mathematics is written, repaired or re-derived: the statement and the proof are the article's own lines, in the article's own notation, and no retained line is substituted or re-flowed.  Retained uncounted as the source-local context the proof needs: lines 746--751; lines 761--772; lines 806--811; lines 836--838.  Nothing was omitted from any retained span, so the package carries no omission record.  No retained line contains a citation, so the wrapper carries no bibliography and no bibliography entry.  No retained line contains a figure environment, an image-inclusion command or drawing code, so no image is delivered and the package declares no asset dependency.  Editorial formatting only, in addition: an AMS \texttt{amsart} preamble that restates the article's own declarations and macros used by the retained lines, namely the environments \texttt{lemma} on separate counters, together with the standard packages those lines need.  The retained environments print in this document as Lemma 1 in order of appearance, whereas the article prints the same items with the numbering of its own full document.  The complete native source of the article as distributed in the arXiv e-print for this exact version is bundled unchanged as \texttt{sources/206251/source0.tex}.
	\begin{equation}\label{GL-oscRealization}
		J_{BA}=\phi(\mathrm{J}_{BA})=\left(\begin{array}{cc}
			\mmu_1-\OAD\OA&\OAD\left((\mmu_1-\mmu_2+1)\ID-\OA\OAD\right)\\
			\OA&(\mmu_2-1)\ID+\OA\OAD
		\end{array}\right) _{AB}
	\end{equation}

	\begin{equation}\label{eq:laxfac}
		L(x)=\left(\begin{array}{cc}
			1&-\OAD\\
			0&\ID
		\end{array}\right)\left(\begin{array}{cc}
			x+\mmu_1&0\\
			\OA&(x+\mmu_2-1)\ID
		\end{array}\right)\left(\begin{array}{cc}
			1&\OAD\\
			0&\ID
		\end{array}\right)\,,
	\end{equation}

		\begin{equation}\label{L1-u1u2}
			L_{1}(x_{1},x_{2})= \left(\begin{array}{cc}
			x_{1}-\OAD^{[1]}\OA^{[1]}&	\OAD^{[1]}\left((x_{1}-x_{2}+1)\ID-\OA^{[1]}\OAD^{[1]}\right)\\
				\OA^{[1]}&(x_{2}-1)\ID+\OA^{[1]}\OAD^{[1]}
			\end{array}\right) \,,
		\end{equation}

		\begin{equation}\label{eq-Rm}
			{\abrOSCR}^-_{12}(x_1,x_2|y_1) L_1(x_1,x_2)L_2(y_1,y_2)=L_1(y_1,x_2)L_2(x_1,y_2){\abrOSCR}^-_{12}(x_1,x_2|y_1)\,,
		\end{equation} 

\begin{lemma}\label{Lemma-equivalence}
	For a shift-invariant operator $\abrOSCR^{-}(x_1,x_2|y_1)=\abrOSCR^{-}(x_1+\lambda,x_2+\lambda|y_1+\lambda)$, the equation in~\eqref{eq-Rm} is equivalent to
	\begin{equation}\label{goal-2a}
		\begin{split}
			&\abrOSCR_{12}^{-}(x_1,x_2|y_1)\left[L_{1}(x_{1},x_{2})+L_{2}(y_{1},y_{2})\right]=\left[L_{1}(y_{1},x_{2})+L_{2}(x_{1},y_{2})\right]\abrOSCR_{12}^{-}(x_1,x_2|y_1),\\
			\end{split}
	\end{equation}
	and
	\begin{equation}\label{goal-2b}
		\begin{split}
			&\abrOSCR_{12}^{-}(x_1,x_2|y_1)\oad_{a}^{[2]}=\oad^{[2]}_a\abrOSCR_{12}^{-}(x_1,x_2|y_1)\qquad \forall a\in \{1,2,\ldots,\M\}\,.
		\end{split}
	\end{equation}
\end{lemma}

\begin{proof}
    
 First, we show that~\eqref{eq-Rm} implies~\eqref{goal-2a} and~\eqref{goal-2b}.
In the following we often suppress the arguments of $\abrOSCR^{-}(x_{1},x_{2}|y_{1})$ and just write $\abrOSCR^{-}$ to shorten the presentation.
Let us consider the following shift of spectral parameters in~\eqref{eq-Rm}:
\begin{equation}\label{variable-shift}
	\begin{split}
		x_{1}\to x_{1}+\lambda,\quad x_{2}\to x_{2}+\lambda,\quad y_{1}\to y_{1}+\lambda \quad \text{and}\quad y_{2}\to y_{2}+\mu\,.
	\end{split}
\end{equation}
Under the shift the Lax matrices  transform as
\begin{equation}
	\begin{split}
		&L(x_{1}+\lambda,x_{2}+\lambda)= L(x_{1},x_{2})+\lambda\ID,\;\quad\qquad\qquad\qquad L(y_{1}+\lambda,x_{2}+\lambda)= L(y_{1},x_{2})+\lambda \ID\\
		&L(y_{1}+\lambda,y_{2}+\mu)=L(y_{1},y_{2})+\lambda\ID+(\mu-\lambda)F,\qquad 	L(x_{1}+\lambda,y_{2}+\mu)= L(x_{1},y_{2})+\lambda\ID+(\mu-\lambda)F
	\end{split}
\end{equation}
where
						\begin{eqnarray}
	F=\begin{pmatrix}
		0&-\OAD\\
		0&\ID
	\end{pmatrix}\,.
\end{eqnarray}
Therefore, using the assumption that $\abrOSCR^{-}$ is shift invariant, the relation~\eqref{eq-Rm} becomes
\begin{equation}\label{consequence-shift-1}
	\begin{split}
		&\abrOSCR_{12}^{-}(x_1,x_2|y_1)\left[\left(L_{1}(x_{1},x_{2})+\lambda\ID\right)\left(L_{2}(y_{1},y_{2})+\lambda\ID+(\mu-\lambda)F^{[2]}\right)\right]\\
		=&
		\left[\left(L_{1}(y_{1},x_{2})+\lambda\ID\right)\left(L_{2}(x_{1},y_{2})+\lambda\ID+(\mu-\lambda)F^{[2]}\right)\right]\abrOSCR_{12}^{-}(x_1,x_2|y_1)\,.
	\end{split}
\end{equation}
By equating the powers of $\lambda$ and $\mu$ one obtains the three conditions
	\begin{align}
			&\abrOSCR_{12}^{-}(x_1,x_2|y_1)\left[L_{1}(x_{1},x_{2})+L_{2}(y_{1},y_{2})\right]=\left[L_{1}(y_{1},x_{2})+L_{2}(x_{1},y_{2})\right]\abrOSCR_{12}^{-}(x_1,x_2|y_1)\label{list-1-first}\\
			&\abrOSCR_{12}^{-}(x_1,x_2|y_1)F^{[2]}=F^{[2]}\abrOSCR_{12}^{-}(x_1,x_2|y_1)\label{list-1-second}\\
			&\abrOSCR_{12}^{-}(x_1,x_2|y_1)L_{1}(x_{1},x_{2})F^{[2]}=L_{1}(y_{1},x_{2})F^{[2]}\abrOSCR_{12}^{-}(x_1,x_2|y_1)\label{list-1-third}\,.
		\end{align}
  Notice that~\eqref{list-1-third} is implied by~\eqref{list-1-first} and~\eqref{list-1-second}. This can be seen as follows.
   First use~\eqref{list-1-second} to bring~\eqref{list-1-third} to the form
   \begin{equation}
       \left[\abrOSCR^{-}(x_1,x_2|y_1)L_{1}(x_{1},x_{2})-L_{1}(y_{1},x_{2})\abrOSCR^{-}(x_1,x_2|y_1)\right]F^{[2]}=0\,.
   \end{equation}
   Next apply~\eqref{list-1-first} and again~\eqref{list-1-second} to obtain
 \begin{equation}
     \abrOSCR^{-}(x_1,x_2|y_1)L_{2}(x_{1},y_{2})F^{[2]}=L_{2}(y_{1},y_{2})F^{[2]}\abrOSCR^{-}(x_1,x_2|y_1)\,.
   \end{equation}
Finally  noting that
\begin{equation}
    L_{2}(y_{1},y_{2})F^{[2]}=(y_{2}-1)F^{[2]}\,,
\end{equation}
we see that~\eqref{list-1-third} follows from~\eqref{list-1-first} and~\eqref{list-1-second}.



   									   
                                                  	


					
   
Now, we show that~\eqref{goal-2a} and~\eqref{goal-2b} imply~\eqref{eq-Rm}.  The logic of this proof consists of three steps. First, one derives an equation equivalent to~\eqref{eq-Rm}. Second,  using the same shift of spectral parameters performed in~\eqref{variable-shift}, one obtains a set of equations that is equivalent to~\eqref{goal-2a} and~\eqref{goal-2b}. Third, one shows that this last set of equations implies~\eqref{eq-Rm}.

\paragraph{Step (i).} One shows that~\eqref{eq-Rm} is equivalent to
\begin{equation}\label{equivalent-def-2}
	\begin{split}
		\mathbf{r}L_{1}(x_{1},x_{2})\left(\begin{array}{cc}
			y_1&0\\
			\OA^{[2]}-\OA^{[1]}&\ID
		\end{array}\right)=
		L_{1}(y_{1},x_{2})\left(\begin{array}{cc}
			x_1&0\\
			\OA^{[2]}-\OA^{[1]}&\ID
		\end{array}\right)\mathbf{r}\,,
	\end{split}
\end{equation}
where
\begin{equation}\label{r-definition}
	\mathbf{r}:=\exp{\left(\OAD^{[2]}\OA^{[1]}\right)}\abrOSCR^{-}\exp{\left(-\OAD^{[2]}\OA^{[1]}\right)}\,.
\end{equation}
This is obtained as follows. The Lax operator is $gl(\M+1)$ invariant, i.e.
\begin{equation}\label{glinv}
	\left(e_{A,B}+ J_{A,B}\right)L(x)=L(x)(e_{A,B}+ J_{A,B})\qquad \forall A,B\in \{0,1,\ldots, \M\}\,.
\end{equation}
Here, $e_{A,B}$ denotes the algebra generators of the fundamental representation with $(e_{AB})_{CD}=\delta_{AC}\delta_{BD}$ and $J_{A,B}$ the algebra generators of the oscillator representation~\eqref{GL-oscRealization}. Consider $\mathbf{t}=(t_{1},\ldots,t_{M})$ such that $\mathbf{t}\OA=\sum_{a=1}^{\M}t_{a}\mathbf{a}_{a}$. By fixing $A=0$ in~\eqref{glinv} such that $J_{0,b}=\mathbf{a}_{b}$, $b=1,\ldots,\M$ and multiplying by $\mathbf{t}$ one finds that
\begin{equation}
	\begin{split}
		\left[\left(\sum_{b=1}^{\M}t_{b}e_{0,b}\right)+ \mathbf{t}\OA\right]L(x)=L(x)\left[\left(\sum_{b=1}^{\M}t_{b}e_{0,b}\right)+ \mathbf{t}\OA\right]\,.
	\end{split}
\end{equation}
This implies that
\begin{equation}
	\begin{split}G^{-1}(\mathbf{t})\exp{\left(\mathbf{t}\OA\right)}L(u)=L(u)G^{-1}(\mathbf{t})\exp{\left(\mathbf{t}\OA\right)}\,.
	\end{split}
\end{equation}
where
\begin{equation}
	G^{-1}(\mathbf{t})=\exp{\left[\sum_{b=1}^{\M}t_{b}e_{0,b}\right]}=\ID+\sum_{b=1}^{\M}t_{b}e_{0,b}
	 \,.
\end{equation}
Then, it follows that
\begin{align}\label{useful-relation}
	G^{-1}(\mathbf{t})L(u)G(\mathbf{t})=\exp{\left(-\mathbf{t}\OA\right)} L(u)\exp{\left(\mathbf{t}\OA\right)}\,.
\end{align}
Equation~\eqref{eq-Rm} is rewritten using~\eqref{eq:laxfac} and~\eqref{useful-relation}, where $t_{a}={\oad}_{a}^{[2]}$ for all $a\in \{1,2,\ldots,\M\}$. One obtains
\begin{equation}
	\begin{split}
		&\abrOSCR^{-}L_{1}(x_{1},x_{2})G(\OAD^{[2]})\left(\begin{array}{cc}
			y_{1}&0\\
			\OA^{[2]}&(y_{2}-1)\ID
		\end{array}\right)G^{-1}(\OAD^{[2]})
		\\=&
		L_{1}(y_{1},x_{2})G(\OAD^{[2]})\left(\begin{array}{cc}
			x_1&0\\
			\OA^{[2]}&(y_{2}-1)\ID
		\end{array}\right)G^{-1}(\OAD^{[2]})\abrOSCR^{-}\,,
	\end{split}
\end{equation}
which is equivalent to
\begin{equation}
	\begin{split}
		&G^{-1}(\OAD^{[2]})\abrOSCR^{-}L_{1}(x_{1},x_{2})G(\OAD^{[2]})\left(\begin{array}{cc}
			y_{1}&0\\
			\OA^{[2]}&(y_{2}-1)\ID
		\end{array}\right)
		\\=&
		G^{-1}(\OAD^{[2]})L_{1}(y_{1},x_{2})G(\OAD^{[2]})\left(\begin{array}{cc}
			x_1&0\\
			\OA^{[2]}&(y_{2}-1)\ID
		\end{array}\right)G^{-1}(\OAD^{[2]})\abrOSCR^{-}G(\OAD^{[2]})\,.
	\end{split}
\end{equation}
Since it has been shown that $\abrOSCR^{-}{\oad}_{a}^{[2]}={\oad}_{a}^{[2]}\abrOSCR^{-}$  for any solution of~\eqref{eq-Rm} one obtains
\begin{equation}
	\begin{split}
		&\abrOSCR^{-}G^{-1}(\OAD^{[2]})L_{1}(x_{1},x_{2})G(\OAD^{[2]})\left(\begin{array}{cc}
			y_{1}&0\\
			\OA^{[2]}&(y_{2}-1)\ID
		\end{array}\right)
		\\=&
		G^{-1}(\OAD^{[2]})L_{1}(x_{1},y_{2})G(\OAD^{[2]})\left(\begin{array}{cc}
			x_1&0\\
			\OA^{[2]}&(y_{2}-1)\ID
		\end{array}\right)\abrOSCR^{-}\,.
	\end{split}
\end{equation}
Next, using~\eqref{useful-relation} one finds
\begin{equation}
	\begin{split}
		&\abrOSCR^{-}\exp{\left(-\OAD^{[2]}\OA^{[1]}\right)}L_{1}(x_{1},x_{2})\exp{\left(\OAD^{[2]}\OA^{[1]}\right)}\left(\begin{array}{cc}
			y_{1}&0\\
			\OA^{[2]}&(y_{2}-1)\ID
		\end{array}\right)
		\\=&
		\exp{\left(-\OAD^{[2]} \OA^{[1]}\right)}L_{1}(y_{1},x_{2})\exp{\left(\OAD^{[2]}\OA^{[1]}\right)}\left(\begin{array}{cc}
			x_1&0\\
			\OA^{[2]}&(y_{2}-1)\ID
		\end{array}\right)\abrOSCR^{-}\,.
	\end{split}
\end{equation}
By computing the commutator
\begin{equation}
	\left[\OAD^{[2]}\OA^{[1]},\mathbf{a}_{b}^{[2]}\right]=-\mathbf{a}_{b}^{[1]}
\end{equation}
 and it follows, by the Hadamard formula,
 that
\begin{equation}
	\exp{\left(\OAD^{[2]}\OA^{[1]}\right)}\mathbf{a}_{a}^{[2]}=\left(\mathbf{a}_{a}^{[2]}-\mathbf{a}_{a}^{[1]}\right)	\exp{\left(\OAD^{[2]}\OA^{[1]}\right)} \qquad \forall a\in \{1,2,\ldots,\M\}\,.
\end{equation}
Therefore, one obtains that
\begin{equation}
	\begin{split}
		&\abrOSCR^{-}\exp{\left(-\OAD^{[2]}\OA^{[1]}\right)}L_{1}(x_{1},x_{2})\left(\begin{array}{cc}
			y_1&0\\
			\OA^{[2]}-\OA^{[1]}&(y_{2}-1)\ID
		\end{array}\right)\exp{\left(\OAD^{[2]}\OA^{[1]}\right)}
		\\=&
		\exp{\left(-\OAD^{[2]}\OA^{[1]}\right)}L_{1}(y_{1},x_{2})\left(\begin{array}{cc}
			x_1&0\\
			\OA^{[2]}-\OA^{[1]}&(y_{2}-1)\ID
		\end{array}\right)\exp{\left(\OAD^{[2]}\OA^{[1]}\right)}\abrOSCR^{-}\,,
	\end{split}
\end{equation}
which can be rewritten as
\begin{equation}
	\begin{split}
		&\exp{\left(\OAD^{[2]}\OA^{[1]}\right)}\abrOSCR^{-}\exp{\left(-\OAD^{[2]}\OA^{[1]}\right)}L_{1}(x_{1},x_{2})\left(\begin{array}{cc}
			y_1&0\\
			\OA^{[2]}-\OA^{[1]}&(y_{2}-1)\ID
		\end{array}\right)
		\\=&
		L_{1}(y_{1},x_{2})\left(\begin{array}{cc}
			x_1&0\\
			\OA^{[2]}-\OA^{[1]}&(y_{2}-1)\ID
		\end{array}\right)\exp{\left(\OAD^{[2]}\OA^{[1]}\right)}\abrOSCR^{-}\exp{\left(-\OAD^{[2]}\OA^{[1]}\right)}\,.
	\end{split}
\end{equation}
By recalling~\eqref{r-definition}
one writes~\eqref{eq-Rm} as
\begin{equation}\label{equivalent-def}
	\begin{split}
		\mathbf{r}L_{1}(x_{1},x_{2})\left(\begin{array}{cc}
			y_1&0\\
			\OA^{[2]}-\OA^{[1]}&(y_{2}-1)\ID
		\end{array}\right)=
		L_{1}(y_{1},x_{2})\left(\begin{array}{cc}
			x_1&0\\
			\OA^{[2]}-\OA^{[1]}&(y_{2}-1)\ID
		\end{array}\right)\mathbf{r}\,.
	\end{split}
\end{equation}
Finally,  observing that
\begin{equation}
	\left(\begin{array}{cc}
			y_1&0\\
			\OA^{[2]}-\OA^{[1]}&(y_{2}-1)\ID
		\end{array}\right)=\left(\begin{array}{cc}
		y_1&0\\
		\OA^{[2]}-\OA^{[1]}&\ID
	\end{array}\right)\left(\begin{array}{cc}
		1&0\\
		0&(y_{2}-1)\ID
	\end{array}\right)\,,
\end{equation}
\eqref{equivalent-def-2} is obtained.

\paragraph{Step (ii).} Next, a set of equations equivalent to~\eqref{goal-2a} and~\eqref{goal-2b} is derived. In particular, it  will be shown that~\eqref{goal-2a} and~\eqref{goal-2b} are equivalent to the following system of equations:
\begin{align}
				&\mathbf{r}{\oad}_{a}^{[1]}\mathbf{a}_{b}^{[1]}={\oad}_{a}^{[1]}\mathbf{a}_{b}^{[1]}\mathbf{r}\qquad \forall a,b \in \{1,\ldots,\M\}
    \label{set1} \\
		&\mathbf{r}\mathbf{a}_{a}^{[2]}=\mathbf{a}_{a}^{[2]}\mathbf{r}\qquad \forall a \in \{1,\ldots,\M\}
    \label{set2} \\
		&\mathbf{r}{\oad}_{a}^{[1]}(x_{2}-x_{1}+\OAD^{[1]}\OA^{[1]})={\oad}_{a}^{[1]}(x_{2}-y_{1}+\OAD^{[1]}\OA^{[1]})\mathbf{r}\qquad \forall a\in \{1,\ldots,\M\}\,.
    \label{set3}
        \end{align}

First, the same shift of parameters~\eqref{variable-shift} in~\eqref{equivalent-def} is applied, obtaining that \begin{equation}
	\begin{split}
		&\mathbf{r}\left(L_{1}(x_{1},x_{2})+\lambda \ID\right) \left[\left(\begin{array}{cc}
			y_1&0\\
			\OA^{[2]}-\OA^{[1]}&(y_{2}-1)\ID
		\end{array}\right)+\lambda\ID+(\mu-\lambda)\begin{pmatrix}
			0&0\\
			0&\ID
		\end{pmatrix}\right]\\
		=&
		\left(L_{1}(y_{1},x_{2})+\lambda \ID\right) \left[\left(\begin{array}{cc}
			x_1&0\\
			\OA^{[2]}-\OA^{[1]}&(y_{2}-1)\ID
		\end{array}\right)+\lambda\ID+(\mu-\lambda)\begin{pmatrix}
			0&0\\
			0&\ID
		\end{pmatrix}\right]\mathbf{r}\,.
	\end{split}
\end{equation}
Then, by equating the powers of $\lambda$ and $\mu$ one obtains
\begin{equation}\label{eqqq}
	\begin{split}
		&\mathbf{r}\left[L_{1}(x_{1},x_{2})+\left(\begin{array}{cc}
			y_1&0\\
			\OA^{[2]}-\OA^{[1]}&(y_{2}-1)\ID
		\end{array}\right)\right]=\left[L_{1}(y_{1},x_{2})+\left(\begin{array}{cc}
			x_1&0\\
			\OA^{[2]}-\OA^{[1]}&(y_{2}-1)\ID
		\end{array}\right)\right]\mathbf{r}
	\end{split}
\end{equation}
and
\begin{equation}\label{eq:rLI}
	\begin{split}
		\mathbf{r}L_{1}(x_{1},x_{2})\begin{pmatrix}
			0&0\\
			0&\ID
		\end{pmatrix}=L_{1}(y_{1},x_{2})\begin{pmatrix}
			0&0\\
			0&\ID
		\end{pmatrix}\mathbf{r}\,.
	\end{split}
\end{equation}
The second equation above is implied by the first one.
By inserting the explicit form of $L_{1}(x_{1},x_{2})$ and $L_{1}(y_{1},x_{2})$, cf.~\eqref{L1-u1u2}, into~\eqref{eqqq} the following equation is obtained
\begin{equation}\label{intermediate}
	\begin{split}
		&\mathbf{r} \left(\begin{array}{cc}
			x_{1}+y_{1}-\OAD^{[1]}\OA^{[1]}&-\OAD^{[1]}\left((x_{2}-x_{1}-1)\ID+\OA^{[1]}\OAD^{[1]}\right)\\
			\OA^{[2]}&(x_{2}+y_{2}-2)\ID+\OA^{[1]}\OAD^{[1]}
		\end{array}\right)
		\\=&
		\left(\begin{array}{cc}
			x_{1}+y_{1}-\OAD^{[1]}\OA^{[1]}&-\OAD^{[1]}\left((x_{2}-y_{1}-1)\ID+\OA^{[1]}\OAD^{[1]}\right)\\
			\OA^{[2]}&(x_{2}+y_{2}-2)\ID+\OA^{[1]}\OAD^{[1]}
		\end{array}\right)\mathbf{r}\,,
	\end{split}
\end{equation}
The elements of the matrix yield the relations in~\eqref{set1}-\eqref{set3}.
Thus, one concludes that they are equivalent to~\eqref{goal-2a} and~\eqref{goal-2b}. 
\paragraph{Step (iii).}  Finally, it is shown that the relations in~\eqref{set1}-~\eqref{set3} imply~\eqref{equivalent-def-2}.
                                  To achieve this, by using~\eqref{eq:rLI} and~\eqref{set2}, equation~\eqref{equivalent-def-2} is reduced to
\begin{equation}\label{eq:relRLv}
	\mathbf{r}L_{1}(x_{1},x_{2})\begin{pmatrix}
		y_{1}\\
		-\OA^{[1]}
	\end{pmatrix}=L_{1}(y_{1},x_{2})\begin{pmatrix}
		x_{1}\\
		-\OA^{[1]}
	\end{pmatrix}\mathbf{r}\,.
\end{equation}
By substituting the matrices $L_{1}(x_{1},x_{2})$ and $L_{1}(y_{1},x_{2})$ one finds
\begin{equation}\label{quasi-final}
	\begin{split}
		&\mathbf{r}
		\begin{pmatrix}
			\OAD^{[1]}\OA^{[1]}\left(x_{2}-y_{1}-x_{1}-1+\OAD^{[1]}\OA^{[1]}\right)\\
				-\OA^{[1]}\left(x_{2}-1-y_{1}+\OAD^{[1]}\OA^{[1]}\right)
            		\end{pmatrix}=
				\begin{pmatrix}
			\OAD^{[1]}\OA^{[1]}\left(x_{2}-y_{1}-x_{1}-1+\OAD^{[1]}\OA^{[1]}\right)\\
			-\OA^{[1]}\left(x_{2}-1-x_{1}+\OAD^{[1]}\OA^{[1]}\right)
            		\end{pmatrix}\mathbf{r}
	\end{split}
\end{equation}
The relation in the upper
component of~\eqref{quasi-final} follows directly from~\eqref{set1}.





Now, it will be shown that the first equation of~\eqref{quasi-final} follows from~\eqref{set3} and~\eqref{set1}. From~\eqref{set3} it follows that
\begin{equation}
    	\mathbf{r}{\OAD}^{[1]}(x_{2}-x_{1}+\OAD^{[1]}\OA^{[1]})\OA^{[1]}={\OAD}^{[1]}(x_{2}-y_{1}+\OAD^{[1]}\OA^{[1]})\mathbf{r}\OA^{[1]}\,,
\end{equation}
by using~\eqref{set1} one obtains
\begin{equation}
    	\mathbf{r}{\OAD}^{[1]}(x_{2}-x_{1}+\OAD^{[1]}\OA^{[1]})\OA^{[1]}={\OAD}^{[1]}\mathbf{r}(x_{2}-y_{1}+\OAD^{[1]}\OA^{[1]})\OA^{[1]}\,.
\end{equation}
And thus
\begin{equation}
    	\mathbf{r}{\OAD}^{[1]}\OA^{[1]}(x_{2}-x_{1}-1+\OAD^{[1]}\OA^{[1]})={\OAD}^{[1]}\mathbf{r}\OA^{[1]}(x_{2}-y_{1}-1+\OAD^{[1]}\OA^{[1]})\,.
\end{equation}
Then, using~\eqref{set1} again, one obtains
\begin{equation}
    	{\OAD}^{[1]}\left\{\OA^{[1]}(x_{2}-x_{1}-1+\OAD^{[1]}\OA^{[1]})\mathbf{r}-\mathbf{r}\OA^{[1]}(x_{2}-y_{1}-1+\OAD^{[1]}\OA^{[1]})\right\}=0\,,
\end{equation}
that implies the first equation of~\eqref{quasi-final}.
\end{proof}

\end{document}
